\chapter{Lo que se verificó y lo que falta}

\begin{enpocas}
Los problemas que aparecieron al llevar el diseño a la computadora, cómo se resolvieron, y lo que queda
pendiente.
\end{enpocas}

\section{Diferencias con el diseño original}

El diseño original del que parte el proyecto tenía nueve detalles que no funcionaban tal cual. Todos
están corregidos en el paquete:
\begin{enumerate}
  \item La etiqueta \cod{v1.17.0} de PX4 no existe; se usa la rama \cod{release/1.17}.
  \item El sensor de la cámara se llama \cod{imager}, no \cod{camera}.
  \item \cod{VehicleStatus} es versionado: el tópico es \cod{/fmu/out/vehicle_status_v1}.
  \item En \cod{x500_mono_cam_down} la cámara está 10~cm \emph{arriba} del centro y ve el propio
        vehículo; aquí va 25~cm abajo.
  \item La configuración de Gazebo de PX4 carga \cod{GstCameraSystem}, que acumula cuadros si el
        codificador se atora: 110~MB/s y el sistema se queda sin memoria en minutos. Aquí Gazebo
        arranca sin él.
  \item \cod{MC_NN_THRST_COEF} trabaja con revoluciones por segundo, no rpm, y faltaban
        \cod{MC_NN_MIN_RPM} y \cod{THR_MDL_FAC}; los valores correctos están en el \emph{airframe}
        4099 y los verifican las pruebas.
  \item La masa del x500 completo es 2.114~kg, no 2.0.
  \item \cod{exportar_tflite.py} usaba \cod{--salida} como nombre del símbolo en C; el módulo exige
        \cod{control_net_tflite}.
  \item PX4 marca sus mensajes con el reloj de pared, no con el de Gazebo; el localizador estima el
        desfase con el mínimo de (llegada $-$ marca).
\end{enumerate}

\section{Lo que se comprobó}

\begin{itemize}
  \item Los programas de modelo, localización, imitación, PPO y exportación reproducen las tablas y
        salidas de referencia (el TFLite difiere de numpy en $1.1\times10^{-6}$).
  \item El detector ajustado: AP50 0.954 en validación; el YOLO de COCO no encontraba a nadie.
  \item La red mejorada: sesgo estacionario de 25.9 a 2.1~cm y misión en 211~s en vez de 427~s.
  \item Las coordenadas GPS de los avisos coinciden con el GPS de PX4, el \emph{home} de ArduPilot y el
        mapa de MAVProxy (prueba \cod{test_origen_gps_igual_al_del_mundo}).
  \item \cod{colcon test}: las 8 pruebas del paquete pasan.
\end{itemize}

\section{Pendientes}

\begin{itemize}
  \item \textbf{Falsas alarmas persistentes.} Grupos de robles que el detector confunde en muchos
        cuadros. Reentrenar con esos ``negativos difíciles'' es el siguiente paso.
  \item \textbf{PPO.} Ya no diverge, pero no mejora a la imitación; falta una recompensa que premie lo
        que la imitación no sabe hacer (por ejemplo, rechazar viento).
  \item \textbf{Seguimiento de cámara en la GUI.} \cod{/gui/follow} no funciona con esta interfaz.
  \item \textbf{Campaña completa.} Con 5 semillas se ven tendencias; las 30 semillas del diseño dan la
        potencia estadística para P1.
\end{itemize}
